% Part: normal-modal-logic
% Chapter: frame-definability
% Section: introduction

\documentclass[../../../include/open-logic-section]{subfiles}

\begin{document}

\olfileid{nml}{frd}{int}

\olsection{مقدمة}

نبحث الآن ما بين علاقة الوصول وصدق !!{formula}s معينة في النماذج
المبنية عليها. ولنبدأ بالانعكاسية: إذا كان لكل~$\OLId{latin-ordinary}{w} \in \OLId{latin-looped}{W}$ أن~$Rww$،
كان كل عالم يصل إلى نفسه. وحين تصدق~$\Box !A$ في~$\OLId{latin-ordinary}{w}$، يكون~$\OLId{latin-ordinary}{w}$
نفسه من العوالم التي يلزم صدق~$!A$ فيها. فانعكاسية~$\OLId{latin-looped}{R}$ في
$\mModel{M}$ تضمن صدق~$\Box !A \lif !A$ أيًا كان العالم~$\OLId{latin-ordinary}{w}$
والصيغة~$!A$. وهذا هو صدق المخطط~$\Box \OLId{latin-ordinary}{p} \lif \OLId{latin-ordinary}{p}$ بجميع حالات
إحلاله في~$\mModel{M}$.

وأما الرجوع من صدق صيغة واحدة إلى الانعكاسية فلا يصح. فقد تصدق
$\Box \OLId{latin-ordinary}{p} \lif \OLId{latin-ordinary}{p}$ في~$\mModel{M}$ ولا تكون~$\OLId{latin-looped}{R}$ انعكاسية.
وشاهده نموذج غير انعكاسي~$\mModel{M}$ تصدق فيه~$\Box \OLId{latin-ordinary}{p} \lif \OLId{latin-ordinary}{p}$
في عوالمه جميعًا: نجعله ذا عالم واحد~$\OLId{latin-ordinary}{w}$ لا يصل إلى نفسه،
ونجعل~$\OLId{latin-ordinary}{w} \in \OLId{latin-looped}{V}(\OLId{latin-ordinary}{p})$. فاختيار قيمة~$\OLId{latin-ordinary}{p}$ على هذا الوجه يكفي لصدق
$\Box \OLId{latin-ordinary}{p} \lif \OLId{latin-ordinary}{p}$، وإن انتفت الانعكاسية.

ولرفع هذا العارض نجرّد النظر من خصوص التعيين~$\OLId{latin-looped}{V}$.
فإذا طلبنا صدق~$\Box \OLId{latin-ordinary}{p} \lif \OLId{latin-ordinary}{p}$ في عوالم~$\mModel{M}$ كلها،
أيًا كانت العوالم الداخلة في~$\OLId{latin-looped}{V}(\OLId{latin-ordinary}{p})$، وجبت انعكاسية~$\OLId{latin-looped}{R}$.
وبيانه أن عدم الانعكاسية يعطي عالمًا~$\OLId{latin-ordinary}{w}$ لا يصدق عنده~$Rww$.
فنجعل~$\OLId{latin-looped}{V}(\OLId{latin-ordinary}{p}) = \OLId{latin-looped}{W} \setminus \{\OLId{latin-ordinary}{w}\}$؛ وحينئذ تصدق~$\OLId{latin-ordinary}{p}$ في غير
$\OLId{latin-ordinary}{w}$، ومن ثم في كل عالم يصل إليه~$\OLId{latin-ordinary}{w}$. إذ يمتنع هنا أن يصل~$\OLId{latin-ordinary}{w}$
إلى~$\OLId{latin-ordinary}{w}$، و$\OLId{latin-ordinary}{w}$ هو العالم الوحيد الذي تكذب فيه~$\OLId{latin-ordinary}{p}$.
وأما~$\OLId{latin-ordinary}{p}$ فتكذب في~$\OLId{latin-ordinary}{w}$ نفسه؛ ولذلك تكذب~$\Box \OLId{latin-ordinary}{p} \lif \OLId{latin-ordinary}{p}$ في~$\OLId{latin-ordinary}{w}$.

لهذا نفرد بنية النموذج، من غير تعيين، باسم \emph{الإطار}.
فالإطار~$\mModel{F}$ زوج~$\tuple{\OLId{latin-looped}{W}, \OLId{latin-looped}{R}}$ من مجموعة عوالم وعلاقة
وصول. وكل نموذج~$\tuple{\OLId{latin-looped}{W}, \OLId{latin-looped}{R}, \OLId{latin-looped}{V}}$ يقال إنه \emph{مبني على}
الإطار~$\tuple{\OLId{latin-looped}{W}, \OLId{latin-looped}{R}}$. فلكل إطار صنف النماذج المبنية عليه،
ولكل صنف أطر صنف النماذج المبنية على واحد من أطره. ونضع
$\mModel{F} \Entails !A$ للدلالة على \emph{صلاحية} !!a{formula}
في الإطار، ومعناه~$\mSat{\OLId{latin-looped}{M}}{!A}$ لكل~$\mModel{M}$ مبني
على~$\mModel{F}$.

وبهذا الترميز، يمكننا إقامة علاقات مطابقة بين ال!!{formula}s وأصناف
الأطر: فمثلًا،~$\mModel{F} \Entails \Box \OLId{latin-ordinary}{p} \lif \OLId{latin-ordinary}{p}$ إذا وفقط إذا كان
$\mModel{F}$ انعكاسيًا.

\end{document}
